\newpage
\section{Related Works} \label{appendix:prior_work}
\textbf{Masked Diffusion Models.}
Masked Diffusion Models (MDMs) were previously built on continuous-time Markov chains over discrete spaces~\citep{hoogeboom2021argmax}. Subsequently, \citet{austin2021structured} introduced D3PM with several families of transition kernels, followed by SEDD~\citep{lou2023discrete}. These discrete diffusions can also be understood through the lens of discrete flows~\citep{campbell2024generative,gat2024discrete}, which parallel flow matching and stochastic interpolants over continuous spaces~\citep{albergo2022building,albergo2023stochastic,lipman2022flow}. Among transition-kernel designs, the absorbing state (\emph{masking}) kernel has been particularly popular due to its simple formulation and training objective. Subsequent work established a theoretical framework for these models~\citep{zheng2023reparameterized,sahoo2024simple,zheng2024masked}. Follow-up works have scaled MDMs up to billions of parameters, showcasing their potential compared to autoregressive models~\citep{nie2025large,dream2025,gemini2025diffusion,labs2025mercury,gong2025diffucoder}.

\textbf{Inference-time strategies for discrete diffusion models.}
A key feature of MDMs is their inference-time flexibility, highlighted by \citet{zheng2024masked,ou2024your} and aligned with our any-order perspective in Section~\ref{sec:prelim_mdm}. Recall the MDM inference procedure from Section~\ref{sec:prelim_mdm}:
\begin{itemize}[leftmargin=15pt,itemsep=0pt,topsep=-0.25em]
\item[(a)] Choose a subset of masked positions $\mathcal{S}\subseteq\{i\mid \rvx_{t_\ell}^i=\mask\}$ to \coloredhl{xblue}{unmask}.
\item[(b)] For each $i\in\mathcal{S}$, \coloredhl{xblue}{unmask} by sampling a clean token from $f_\theta^i(\cdot\mid \rvx_{t_\ell})\in\Delta(\mathcal{V})$ and setting $\rvx_{t_\ell}^i$ accordingly.
\end{itemize}
The choice of $\mathcal{S}$ is entirely flexible, and an appropriate selection can substantially improve downstream task performance. Since our work focuses on self-correction, i.e., re-masking, we only briefly survey prior work on unmasking strategies \citep{kim2025train,peng2025pathplanningmaskeddiffusion,chang2022maskgit,ben2025accelerated,rout2025anchored}. A complementary line of work studies inference-time planning strategies for discrete diffusion with uniform transition kernels \citep{liu2024think,varma2024glauber}; these methods fall outside the scope of MDMs.

% \rebuttal{
% \textbf{Comparison to \citet{huang2025don}.} For clarity, we provide a comprehensive comparison with \citet{huang2025don}. The overall training pipeline is largely similar but differs in several key aspects: (1) we provide theoretical guarantees for our training method (PRISM), and (2) we use a pretrained MDM to sample from $\rvz$ to $\rvy$, whereas \citet{huang2025don} construct $\rvy$ via a random corruption process. Nevertheless, their large-scale experiments on LLaDA-8B-Instruct cover a broader range of tasks and datasets, while our work focuses on relatively smaller-scale experiments.}

\subsection{Details on prior work on remasking strategy (self-correction) for MDM}
In this section, we detail our discussion in Section~\ref{sec:prior_work_main}, focusing on their different notions of per-token qualities. We recall the notion of per-token quality stated in Section~\ref{sec:prior_work_main}.
\begin{equation*}
    g_\star\colon (\mathcal{V}\cup\{\mask\})^L \to [0,1]^L,\quad \xxgreen{g^i_\star(\rvy) \approx (\text{Per-token quality})}=(\text{``how likely'' $\rvy^i$ is given $\rvy$}).
\end{equation*}
\begin{itemize}[leftmargin=*,itemsep=0pt]
\item \textbf{DFM}~\citep{gat2024discrete} and \textbf{ReMDM}~\citep{wang2025remasking} randomly remask clean tokens in $\rvy$, which is equivalent to assigning i.i.d.\ per-token qualities $g_\star^i(\rvy)\sim \operatorname{Unif}[0,1]$.
\item \textbf{P2-Self}~\citep{peng2025pathplanningmaskeddiffusion} defines per-token quality using the pretrained MDM’s score for the observed token under input $\rvy$, e.g., the probability $f_\theta^i(\rvy^i \mid \rvy)$. Given that $\rvy^i\neq\mask$, the model was never trained on this index in this state; it only encodes a meaningful quantity (unmasking posterior) when $\rvy^i=\mask$. This notion of per-token quality is therefore not grounded.
\item \textbf{ReMDM-conf}~\citep{wang2025remasking} takes the per-token quality to be the likelihood at the step the token was unmasked: if position $i$ was unmasked at $t_k>t_\ell$ (i.e., an earlier step), set $g_\star^i := f_\theta^i(\rvx_{t_k}^i \mid \rvx_{t_k})=f_\theta^i(\rvx_{t_\ell}^i \mid \rvx_{t_k})$. As noted, this ignores the current state $\rvx_{t_\ell}$ and can therefore be misaligned with the model’s present beliefs.
\item \textbf{Token Critic}~\citep{lezama2022improved,lezama2022discrete} and \textbf{P2-train}~\citep{peng2025pathplanningmaskeddiffusion} train an auxiliary scorer for the likelihood of a given prediction of clean sequence $\hat{\rvx}$. Given $\rvx_{t_\ell}$, they first form a full completion $\hat{\rvx}$ by unmasking each mask token at position $i$ from $f_\theta^i(\cdot\,\mid\,\rvx_{t_\ell})$, then evaluate quality by passing this $\hat{\rvx}$ to the external module. A key limitation is that $\hat{\rvx}$ is drawn from per-token unmasking posteriors, not the ground-truth posterior across positions; thus, even with a perfect pretrained MDM, $\hat{\rvs}$ can deviate substantially from the true likelihood $\rvx\sim p_{\mathrm{data}}$, losing a theoretical guarantee.
\item \textbf{GIDD}~\citep{von2025generalized} introduces a family of transition kernels that linearly combine uniform and absorbing (masking) kernels. Models pretrained under this kernel encode posteriors at clean-token positions--however, this requires training under the new kernel and is not directly applicable to an off-the-shelf pretrained MDM, altogether losing the benefit of MDM over other types of kernels.
\item \textbf{Informed corrector}~\citep{zhao2024informed} employs a hollow-transformer $g_{\phi'}$ satisfying $g_{\phi'}^i(\rvy)=g_{\phi'}^i(\rvy\oplus \mask_i)\approx p(\rvx^i=\rvy^i\mid \rvy\oplus \mask_i)$, enabling the same training loss as MDM. This, however, mandates a modified architecture (attention tweak) that processes $\rvy$ as if position $i$ were masked, necessitating retraining.
\end{itemize}
In summary, while these methods improve over vanilla MDM inference without remasking, they either substantially reshape the MDM pipeline or rely on imprecise notions and proxies for per-token quality.